\scanpage{14}

% Wholly handwritten page: the main hand is unnamed.
% The following clipped left-edge traces have no secure lexical identification.
% Their textual status is uncertain; the phantom alternatives do not deny the ink.
\reversemarginpar
\margin{\figure{\unsure{\begin{tikzpicture}[x=.35mm,y=.35mm,line width=.6pt,line cap=round]
\draw (0,16) -- (0,14) -- (6,13) .. controls (9,9) and (9,7) .. (6,7)
 .. controls (3,9) and (1,12) .. (0,14);
\draw (0,4) .. controls (5,0) and (10,-1) .. (11,1)
 .. controls (12,3) and (10,5) .. (7,5);
\draw (1,-3) .. controls (3,-4) and (3,-6) .. (1,-7);
\end{tikzpicture}|\phantom{\rule{4mm}{9mm}}}}}
\normalmarginpar
b) $\forall\,F$ injective on $\unsure{Y}$, $\unsure{q|g}^{*}(F)$ is injective and $F\unsure{\mathord{-}|\to}q_{*}q^{*}F$ is an epimorphism (example: $q$ inclusion of an open $U\hookrightarrow e_{Y}$).
One takes
\[
\mathbb R\Gamma_{Y}\bigl(\operatorname{Ker}\bigl(\underbrace{C(F}_{\text{injective resolution}})\to
q_{*}q^{*}\unsure{(|C}C(F)\bigr)\bigr)\bigr)
\]
% Small detached right-margin trace beside/below the first display; status uncertain.
\margin{\figure{\unsure{\begin{tikzpicture}[x=.35mm,y=.35mm,line width=.35pt,line cap=round]
\draw[gray] (0,5) .. controls (1,4) and (2,5) .. (3,5);
\fill (2,2.5) circle (.35);
\draw[gray] (2,2.2) -- (1.5,.5);
\fill[gray] (1,-1) ellipse (.3 and .65);
\end{tikzpicture}|\phantom{\rule{2mm}{4mm}}}}}

In the general case, the difficulty comes from the fact that the cone of a morphism of complexes (such as
% Clipped left-edge upright, lower bar and point; no secure textual identity.
\reversemarginpar
\margin{\figure{\unsure{\begin{tikzpicture}[x=.35mm,y=.35mm,line width=.6pt,line cap=round]
\draw (0,8) -- (0,2);
\draw (0,0) -- (5,.4);
\fill (0,-5) ellipse (.45 and .7);
\end{tikzpicture}|\phantom{\rule{3mm}{6mm}}}}}
\normalmarginpar
\[
F\to q_{*}\deleted{\unsure{C|(}}\bigl(q^{*}(F)\deleted{\unsure{)|C}}\bigr)
\deleted{\unsure{*|q_{*}}}\qquad\text{)}
\]
is not functorial (in the
% Detached dark left-edge point, textual status uncertain.
\reversemarginpar
\margin{$\unsure{\cdot|\phantom{\cdot}}$}
\normalmarginpar
derived category) with respect to the arrow whose cone one wants to take.
% Detached left-margin point near source y=914; textual status uncertain.
\reversemarginpar
\margin{$\unsure{\cdot|\phantom{\cdot}}$}
\normalmarginpar

% Clipped curved trace at the left edge beside this paragraph.
\reversemarginpar
\margin{\figure{\unsure{\begin{tikzpicture}[x=.4mm,y=.4mm,line width=.6pt,line cap=round]
\draw (0,3) .. controls (4,1) and (4,-2) .. (0,-3);
\end{tikzpicture}|\phantom{\rule{2mm}{3mm}}}}}
\normalmarginpar
And yet, in the present particular case\unsure{\phantom{,}|,} there should be a functorial choice.
% The detached point lies below-left of the hyphen between t and ce.
IS$\underset{\hspace{-.35em}\unsure{\cdot|\phantom{\cdot}}}{\text{-}}$this
\unsure{\.{e}vident|evident|evident} ?
\textsuperscript{\unsure{'|\phantom{'}}}
% A further clipped curved trace lies at the left edge beside the preceding question.
\reversemarginpar
\margin{\figure{\unsure{\begin{tikzpicture}[x=.4mm,y=.4mm,line width=.6pt,line cap=round]
\draw (0,3) .. controls (3,1) and (3,-2) .. (0,-3);
\fill (0,-9) ellipse (.25 and .4);
\end{tikzpicture}|\phantom{\rule{2mm}{5mm}}}}}
\normalmarginpar

% Isolated point in the gap above the rubric, source (510,1240).
\noindent\hspace{.43\linewidth}$\unsure{\cdot|\phantom{\cdot}}$\par
% Two small boundary shapes beside the rubric underline; pen origin uncertain.
\reversemarginpar
\margin{\figure{\unsure{\begin{tikzpicture}[x=.4mm,y=-.4mm,line width=.5pt,line cap=round]
\draw (0,0)--(2,1.2);
\fill (1,12) circle (.8);
\end{tikzpicture}|\phantom{\rule{2mm}{6mm}}}}}
\normalmarginpar
\unsure{Q\uline{uestion for Illusi}e|Q\uline{uestion for Illusie}} : In his \unsure{th\.{e}ory|theory|theory} of deformations of flat \unsure{group|group}
% Clipped left-edge thick curved trace and small lower point.
\reversemarginpar
\margin{\figure{\unsure{\begin{tikzpicture}[x=.4mm,y=.4mm,line width=.8pt,line cap=round]
\draw (0,3) .. controls (2,1) and (2,0) .. (0,-1);
\fill (0,-6) ellipse (.25 and .4);
\end{tikzpicture}|\phantom{\rule{2mm}{4mm}}}}}
\normalmarginpar
schemes, he runs into $H^{3}(B_{\unsure{G|C}}/X,\quad)$ resp. $\operatorname{Ext}^{2}(\unsure{K|X};,\quad)$.
Can one short-circuit his theory via the theory (assumed written) of Gr-stacks - resp. via the theory of Picard stacks ? I have

% Two detached bottom points, below the faint horizontal paper/image rule.
\begin{center}
$\unsure{\bullet\,\cdot|\phantom{\bullet\,\cdot}}$
\end{center}
% The faint horizontal gray rule near the foot is not assigned a textual function.
% archive sheet number: 13
